% Lecture example set: SC4061-L05
% Sources: L3 pp. 32--45; L4 pp. 1--19; checked discussion on 2026-09-08
% Human verified: Yes

\clearpage

\exercisemeta
  {SC4061-L05-EX}
  {2026-09-08}
  {L3 (full) Frequency slides pp.~32--45；L4 (full) Edge slides pp.~1--19；课堂讨论中的自编检验例}
  {四道检验题答案已完成并核对}
  {已确认（2026-09-08）}

\exercisequestion{SC4061-L05-E01}{Local Edge 的 Fourier/Wavelet Representation}

\textbf{类型：}lecture quiz / self-contained check（基于 L3 p.~32）。

\textbf{对应知识点：}
\relatedtopic{SC4061-L05-T01}{Dense Fourier Spectrum 与 Wavelet Localization}

\subsubsection*{题目}

一幅图像只有两种 intensity，但包含一个局部 sharp boundary。解释为什么它在 spatial domain 很简单，而 Fourier spectrum 仍可能很 dense；并说明 wavelet 相比 Fourier 额外保留了什么信息。

\begin{checkedsolution}
Fourier basis functions 是覆盖全图的 global sinusoids。局部 abrupt change 无法由少数平滑 modes 表示，需要许多不同 frequencies 的 Fourier coefficients 共同重建，因此 spectrum 可能很 dense。Wavelet 使用 localized、multi-scale basis，其 coefficients 不仅对应 frequency/scale，还保留 approximate spatial location。

这里应说 ``许多 Fourier coefficients/components''，而不是``许多 spectra''：一幅图像只有一个 spectrum，其中包含许多 coefficients。
\end{checkedsolution}

\textbf{检查点：}DFT 始终可以计算；``可计算''与``该 domain 中表示稀疏、处理方便''不是同一件事。

\exercisequestion{SC4061-L05-E02}{Sobel Gradient 与 Edge Direction}

\textbf{类型：}self-contained check（基于 L4 p.~7）。

\textbf{对应知识点：}
\relatedtopic{SC4061-L05-T05}{Prewitt 与 Sobel Gradient Filters}

\subsubsection*{题目}

对 patch
\[
  P=\begin{bmatrix}
  20&20&20\\
  20&20&20\\
  80&80&80
  \end{bmatrix}
\]
使用讲义中的 Sobel kernels，求 $G_x,G_y$、$|\nabla f|$、gradient direction 和 edge direction。

\begin{checkedsolution}
\[
  G_x=0,
\]
\[
  G_y=(-20-40-20)+(80+160+80)=240.
\]
因此
\[
  |\nabla f|=\sqrt{0^2+240^2}=\boxed{240},\qquad
  \alpha=\operatorname{atan2}(240,0)=\boxed{90^\circ}.
\]
Gradient 沿竖直方向，edge tangent 与 gradient 垂直，所以这是 horizontal edge（水平边缘）。
\end{checkedsolution}

\textbf{检查点：}$G_y$ 响应沿 $y$ 的强度变化，但检测出的 edge 本身沿 $x$ 方向延伸；不要把 gradient direction 与 edge direction 混淆。

\exercisequestion{SC4061-L05-E03}{LoG 的 Smoothing 与 Zero Crossing}

\textbf{类型：}concept check（基于 L4 pp.~9--10）。

\textbf{对应知识点：}
\relatedtopic{SC4061-L05-T06}{Laplacian of Gaussian 与 Zero Crossing}

\subsubsection*{题目}

为什么 LoG 在 second derivative 前引入 Gaussian smoothing？对于 step edge，应使用 response peak 还是 zero crossing 确定 edge location？

\begin{checkedsolution}
Second derivative 会放大 high-frequency noise；Gaussian smoothing 先衰减噪声，再由 Laplacian 提取 curvature。由于
\[
  f*(\nabla^2G_\sigma)=\nabla^2(f*G_\sigma),
\]
两步可合并为一次 LoG convolution。Step edge 的 LoG response 在边缘两侧形成正负 lobes，edge centre 位于两者之间的 \boxed{\text{zero crossing}}，而不是任一 peak。
\end{checkedsolution}

\textbf{检查点：}LoG 的 Gaussian 只降低 sensitivity，不会使二阶导数对噪声完全免疫。

\exercisequestion{SC4061-L05-E04}{Hysteresis Thresholding 的 Strong/Weak Edges}

\textbf{类型：}self-contained check（基于 L4 pp.~17--18）。

\textbf{对应知识点：}
\relatedtopic{SC4061-L05-T07}{Canny Edge Detector}

\subsubsection*{题目}

设 $t_L=40,t_H=100$。Pixel A 的 magnitude 为 $130$；B 为 $70$，并通过 weak-edge chain 与 A 连通；C 为 $70$，但完全孤立；D 为 $20$。Hysteresis thresholding 后哪些 pixels 被保留？

\begin{checkedsolution}
A 满足 $M\ge t_H$，是 strong edge，直接保留。B 位于双阈值之间，但通过 weak-edge chain 与 strong edge A 连通，因此保留。C 虽然也位于双阈值之间，却没有 strong-edge anchor，删除；D 低于 $t_L$，删除。故最终保留
\[
  \boxed{\text{A 和 B}}.
\]
\end{checkedsolution}

\textbf{检查点：}Weak edge 不是一律保留或删除；关键是它是否通过 connectivity 与 strong edge 相连。
